\chapter{文件系统}
\label{chap:fs}

\begin{introduction}[本章内容]
  \item xv6 文件系统的分层
  \item 磁盘布局与 \code{mkfs}
  \item 块缓存、日志与 inode
  \item 块设备：从 ramdisk 到 SD 卡
  \item 根文件系统放在哪里
  \item VFS 中间层：procfs、ext2、FAT32
  \item FAT32 与启动分区：bootfs、BPB、FAT 表、目录项与长文件名
  \item xv6 根文件系统：superblock、inode、目录、路径查找与在线更新
  \item 烧录 SD 卡
\end{introduction}

\section{分层结构}

xv6 的文件系统是一叠职责单一的层（图~\ref{fig:fslayers}）。本项目保留了中间各层不动，在上面加了一层 VFS，在下面把块设备从 VirtIO 磁盘换成了真正的 SD 卡。

\begin{figure}[htbp]
  \centering
  \begin{tikzpicture}[node distance=2mm]
    \tikzset{layer/.style={box, minimum width=78mm, text width=74mm, align=left}}
    \node[layer] (sys) {\textbf{系统调用}\quad \code{open/read/write/mkdir/mount}};
    \node[layer, below=of sys] (file) {\textbf{文件描述符}\quad \code{struct file}：inode、设备、管道、vnode、socket、PTY};
    \node[layer, below=of file, fill=structurecolor!22] (vfs) {\textbf{VFS}\quad 挂载表、\code{vnode\_ops}：procfs、ext2、FAT32};
    \node[layer, below=of vfs] (path) {\textbf{路径名}\quad \code{namei}、目录项};
    \node[layer, below=of path] (inode) {\textbf{inode}\quad \code{ilock/readi/writei}、inode 缓存};
    \node[layer, below=of inode] (log) {\textbf{日志}\quad \code{begin\_op/log\_write/end\_op}};
    \node[layer, below=of log] (bio) {\textbf{块缓存}\quad \code{bread/bwrite}，\code{NBUF}=30};
    \node[layer, below=of bio, fill=structurecolor!22] (map) {\textbf{根块设备}\quad \file{rootdev.c}，\code{rootdev\_rw()}：原始分区 / 裸镜像};
    \node[layer, below=of map, fill=structurecolor!22] (sd) {\textbf{SD 卡驱动}\quad \code{sdsector()}：SDHOST，CMD17/CMD24};
  \end{tikzpicture}
  \caption{文件系统分层（深色为本项目新增或替换的层）}
  \label{fig:fslayers}
\end{figure}

\section{磁盘布局}
\label{sec:layout}

xv6 原生文件系统的块大小是 1\,KiB，磁盘依次为：

\begin{txtcode}
[ 启动块 | superblock | 日志 | inode 块 | 空闲位图 | 数据块 ]
   0          1          2..
\end{txtcode}

superblock 记录各区的起始块号和大小，魔数为 \code{0x10203040}。几个参数与原版不同：

\begin{itemize}
  \item \code{FSSIZE}$=32768$：文件系统 32\,MiB。原来只有 1000 块，加入网络、SSH 等程序后，打包就用掉了 986 块，首次启动创建 \code{/etc} 配置时 panic “balloc: out of blocks”；
  \item \code{DIRSIZ}$=30$：文件名最长 30 字节，目录项从 16 字节变为 32 字节，为了让 TFTP 下载的 \code{kernel8-xv6\_wifi.img} 这样的长文件名能通过普通 \code{ls} 显示；
  \item inode 有 12 个直接块和 1 个一级间接块，单个文件最大 $12+256=268$ 块，即 268\,KiB。
\end{itemize}

\begin{note}
  分区再大也不会改变可分配的块数——它由 superblock 里的 \code{size} 字段决定。修改 \code{FSSIZE} 或 \code{DIRSIZ} 之后必须重新生成并整体烧录 \code{fs.img}，只替换内核会让新内核去解释旧格式的磁盘。
\end{note}

\section{块缓存、日志与 inode}

这三层基本沿用 xv6 原样，这里只说明本项目中与它们相关的两点。

\paragraph{日志.} 一次文件系统操作可能修改多个块（例如创建文件要改 inode、目录、位图）。\code{begin\_op()}/\code{end\_op()} 之间的所有修改先写进日志区，提交时一次性写回原位置；启动时 \code{fsinit()} 检查日志，把已提交但未写完的事务重做一遍。这样即使在写的过程中断电，文件系统也只会处于“操作前”或“操作后”两种状态之一。对 SD 卡来说这很重要：真机经常是直接拔电源关机的。

\paragraph{inode 睡眠锁.} inode 操作可能触发 SD 卡读写，耗时很长，不能持有自旋锁，所以每个内存 inode 有一把睡眠锁。项目中出现过一次典型的锁配对错误：在为文件描述符层加入 \code{FD\_DEVICE}、\code{FD\_VNODE}、\code{FD\_SOCKET} 等分支时，\code{fileread()} 的 \code{FD\_INODE} 分支丢了 \code{iunlock()}。

\begin{ccode}
} else if(f->type == FD_INODE){
  ilock(f->ip);
  if((r = readi(f->ip, 1, addr, f->off, n)) > 0)
    f->off += r;
  iunlock(f->ip);       // 这一行被误删
}
\end{ccode}

现象非常迷惑：启动卡在 “exec-login: inode found” 之后，调试输出显示 \code{/bin/login} 的 inode 被 PID 1 锁着，而 \code{valid=0}。原因是 \code{init} 读 \code{/etc/fstab} 时没有解锁，关闭文件后引用计数归零，这个 inode 缓存槽被 \code{iget()} 复用给了 \code{/bin/login}——槽换了主人，锁却还“记得”上一任的持有者。

\begin{note}
  \code{iput()} 不能代替 \code{iunlock()}：前者管“有多少引用”，后者管“是否正在独占访问内容”。一个有用的调试断言是：\code{iget()} 复用 \code{ref==0} 的槽时，若发现锁仍被持有，立即 panic 并打印旧的 dev/inum/owner。
\end{note}

\section{块设备：从 ramdisk 到 SD 卡}

\subsection{三个阶段}

块设备驱动经历了三次替换：

\begin{enumerate}
  \item \textbf{ramdisk}：刚移植到 \code{raspi3b} 时，用 QEMU 的 loader 把 \code{fs.img} 装到物理地址 \code{0x07000000}，\code{bread/bwrite} 直接 \code{memmove}。简单，但改动不会写回宿主机；
  \item \textbf{Arasan SDHCI}（\file{kernel/sd.c}）：BCM2837 的 EMMC 控制器，GPIO48～53 设为 ALT3。QEMU 用 \code{-drive if=sd} 把 \code{fs.img} 接到 SD 总线，数据终于可以持久保存；
  \item \textbf{SDHOST}（\file{kernel/sdhost.c}）：后来要驱动板载 Wi-Fi，而 Wi-Fi 芯片只能接在 Arasan 控制器上（GPIO34～39）。于是外置 SD 卡改由 Broadcom 自己的 SDHOST 控制器（\code{0x3f202000}）驱动，GPIO48～53 改为 ALT0。
\end{enumerate}

两种 SD 驱动对上都提供同一个接口：

\begin{ccode}
int sdsector(uint32 sector, void *buffer, int write);   // 读写一个 512 字节扇区
\end{ccode}

xv6 的块是 1\,KiB，因此一个块对应两个扇区：块 $N$ 是扇区 $2N$ 和 $2N+1$。

\subsection{SD 卡初始化}

无论哪个控制器，SD 卡协议层的初始化顺序都一样（SD 物理层规范\cite{sdspec}）：

\begin{txtcode}
400 kHz 识别时钟
CMD0                复位，卡进入 idle
CMD8                检查电压范围，确认 SD 2.0
CMD55 + ACMD41      反复发送直到卡完成上电；OCR.CCS=1 表示 SDHC（按扇区寻址）
CMD2                读 CID
CMD3                取得相对地址 RCA
CMD7                选中卡，进入 transfer 状态
(SDSC) CMD16        设置块长 512
切换到 25 MHz 工作时钟
\end{txtcode}

读一个扇区发 CMD17，写一个扇区发 CMD24，数据通过数据寄存器以轮询 PIO 方式搬运 128 个 32 位字。驱动用一把自旋锁串行化对控制器的访问。

\subsection{命令超时与恢复}

真机在 shell 里执行 \code{ls} 时偶尔 panic “fat32: read FS.IMG”。在驱动里加入事务阶段编号后得到：

\begin{txtcode}
sd: I/O failure lba=295956 write=0 phase=1 status=17f0207 irq=18000 control1=e0807
\end{txtcode}

\code{phase=1} 表示失败发生在 CMD17 的命令阶段；中断状态 \code{0x18000} 是 ERROR 加 COMMAND TIMEOUT。更麻烦的是，超时后控制器的命令/数据状态机停在 inhibit 状态，之后的每条命令都会失败。仅仅清中断位不够，还要用 \reg{CONTROL1} 的软件复位位分别复位命令线和数据线，等硬件自动清零这两位。修复后的 \code{sdsector()} 对同一扇区最多重试三次，每次失败后先 \code{recover\_io()}。这只复位主机控制器内部的状态机，卡本身仍处于 transfer 状态，不需要重新初始化。

\section{根文件系统放在哪里}
\label{sec:rootfs}

启动时 \code{rootdev\_init()} 读 SD 卡第 0 扇区，判断根文件系统用哪种布局（表~\ref{tab:rootfs}），此后由 \code{rootdev\_rw()} 把 xv6 块号翻译成 SD 卡的绝对扇区号。这两个函数都在 \file{kernel/rootdev.c} 里。

\paragraph{根块设备不是 FAT32 的一部分.} \file{rootdev.c} 不解释任何文件系统格式。它只读 MBR，用 xv6 superblock 的魔数认出根分区，再把 FAT32 启动分区交给 \file{fat32.c} 挂起来，供 \code{/boot} 和固件加载使用。根分区的块里装的始终是 xv6 自己的 superblock、inode 和目录（\ref{sec:xv6fs} 节），与 FAT32 毫无关系。块缓存的读写入口 \code{rootdev\_rw()} 也只做地址换算，不调用任何 FAT32 函数。

\begin{table}[htbp]
  \centering
  \caption{根文件系统的两种存放方式}
  \label{tab:rootfs}
  \small
  \begin{tabularx}{\linewidth}{@{}lXl@{}}
    \toprule
    方式 & 识别条件与地址换算 & 适用场景 \\
    \midrule
    原始分区 & MBR 中某个非 FAT32 主分区的第 1 块是 xv6 superblock（魔数与大小匹配）；扇区 $=$ 分区起点 $+$ 块号$\times 2$ & 真机 \\
    裸镜像 & 没有 MBR；扇区 $=$ 块号$\times 2$ & QEMU \\
    \bottomrule
  \end{tabularx}
\end{table}

\paragraph{原始分区模式}是真机上的做法：SD 卡分成 4 个区，xv6 文件系统独占一个分区，与 FAT32 启动分区互不干扰。识别时只认 superblock 魔数，而不看 MBR 类型（macOS 只能建 \code{0x83} 类型的分区，与 ext2 相同）。启动日志会显示：

\begin{txtcode}
rootdev: raw xv6 root p2 lba=208896 sectors=819200 size=400 MiB
rootdev: bootfs p1 lba=2048 mounted as FAT32
\end{txtcode}

\section{VFS 中间层}
\label{sec:vfs}

\subsection{动机}

最初的 ext2 驱动只能通过两个专用系统调用 \code{ext2read}、\code{ext2readdir} 访问，普通的 \code{ls}、\code{cat} 完全不知道它的存在。VFS 层把非原生文件系统包装成 \emph{vnode}，让它们通过统一的文件描述符接口工作：

\begin{ccode}
struct vnode_ops {                       // kernel/vfs.c，参数是挂载点内的相对路径
  int (*stat)(char*, struct stat*);
  int (*read)(char*, uint64 off, void*, int n);
  int (*write)(char*, struct fat32_file*, uint64 off, int user_src, uint64 src, int n);
  int (*readdir)(char*, int index, void*);
  int (*create)(char*);
  int (*truncate)(char*, struct fat32_file*);
  int (*rename)(char*, char*);
};
\end{ccode}

这些操作以路径为参数而不是以 inode 为参数，后端每次按路径重新查找。这是一个极简的 VFS：没有 dentry 缓存，也没有把原生 xv6 inode 纳入 vnode，但足以让 \code{ls}、\code{cat}、\code{mv} 等普通命令透明地访问其他文件系统。

\code{struct file} 增加了 \code{FD\_VNODE} 类型，\code{fileread()}、\code{filestat()}、\code{fileclose()} 遇到它就转发给对应的 \code{vnode\_ops}。

\subsection{挂载表与路径路由}

\code{sys\_open()} 先调用 \code{vfsopen()}，它在挂载表里做\emph{最长前缀匹配}：

\begin{txtcode}
open("/boot/config.txt")
  -> findmount(): 绝对路径 /boot/config.txt，匹配挂载点 /boot，相对路径 /config.txt
  -> fat32_ops.stat("/config.txt")
\end{txtcode}

返回值约定是这一层设计的关键：\code{1} 表示 VFS 已打开；\code{-1} 表示路径属于某个挂载点但打开失败（不能退回原生根文件系统，否则会打开挂载点下面被遮住的同名文件）；\code{0} 表示不属于任何挂载点，继续走原生的 \code{namei()}。

\subsection{三个后端}

\begin{itemize}
  \item \textbf{procfs}（\code{/proc}）：没有磁盘数据，打开目录时遍历进程表生成每个 PID 的目录，读 \code{status} 时持锁生成快照。另有 \code{/proc/meminfo}、\code{/proc/iomem}、\code{/proc/devices} 等诊断文件；
  \item \textbf{ext2}（\code{/mnt/ext2}）：支持 1/2/4\,KiB 块，遇到 ext4 extent 等不兼容特性就拒绝挂载，可以和 Linux 交换文件。最初只能读、只支持一级间接块；后来补成了可读写、带外置事务日志的后端，见第~\ref{chap:optional} 章；
  \item \textbf{FAT32}（\code{/boot}）：挂载启动分区，支持根目录的读、创建、截断、顺序写和重命名，支持 ASCII 长文件名。这让系统可以在运行时更新自己的内核：\code{tftp} 下载 \code{kernel8-xv6\_wifi.img} 到 \code{/boot}，重启即生效。
\end{itemize}

\subsection{由用户态决定挂载}

挂载策略仿照嵌入式 Linux：内核只负责探测驱动，\code{init} 读取 \code{/etc/fstab} 决定挂载什么。首次启动时 \code{init} 创建：

\begin{txtcode}
proc    /proc      procfs ro 0 0
bootfs  /boot      fat32  rw 0 0
ext2    /mnt/ext2  ext2   rw 0 0
\end{txtcode}

然后逐行调用 \code{mount(source, target, fstype, flags)}。启动日志出现 “vfs: mounted fat32 at /boot read-write” 才说明 FAT32 真的挂上了；否则 \code{/boot} 只是 \code{init} 创建的空目录。

\section{FAT32 与启动分区}
\label{sec:fat32}

前面几节讲的是 xv6 自己的文件系统。SD 卡上还有一个 xv6 无法回避的 FAT32 分区：树莓派的固件只认 FAT32，内核镜像、\file{config.txt}、Wi-Fi 固件都必须放在这里。本节先看这个分区里放了什么、谁来读，再看 FAT32 在磁盘上的组织，最后逐个讲 \file{kernel/fat32.c} 里的函数。

\subsection{bootfs 里有什么，谁来读}

图~\ref{fig:bootfs} 按读者把启动分区的文件分成三组。

\begin{figure}[htbp]
  \centering
  \begin{tikzpicture}[x=1mm,y=1mm, font=\footnotesize,
      rd/.style={-{Stealth[length=2mm]}, thick, draw=structurecolor!70!black},
      wr/.style={{Stealth[length=2mm]}-{Stealth[length=2mm]}, thick, draw=orange!80!black},
      f/.style={font=\scriptsize\ttfamily, anchor=west, inner sep=1pt}]
    \draw[rounded corners=3pt, draw=gray!60, fill=gray!4] (48,0) rectangle (100,74);
    \node[lab, anchor=north] at (74,73.5) {SD 卡 p1：BOOTFS（FAT32 根目录）};
    % group A: firmware boot
    \draw[fill=gray!10, draw=gray!50] (51,52) rectangle (97,69);
    \node[f] at (52,66.5) {bootcode.bin};
    \node[f] at (52,63) {start.elf  fixup.dat};
    \node[f] at (52,59.5) {config.txt};
    \node[f] at (52,56) {kernel8-xv6\_wifi.img};
    \node[lab, anchor=east] at (96.5,53.8) {启动阶段};
    % group B: drivers
    \draw[fill=structurecolor!8, draw=structurecolor!50] (51,24) rectangle (97,49);
    \node[f] at (52,46) {BCM43430.BIN/.TXT/.CLM};
    \node[f] at (52,42.5) {BCM43455.BIN/.TXT/.CLM};
    \node[f] at (52,39) {MT7601U.BIN};
    \node[lab, anchor=east] at (96.5,26) {驱动固件};
    % group C: runtime
    \draw[fill=orange!8, draw=orange!50] (51,3) rectangle (97,21);
    \node[f] at (52,18) {camera.bmp  camera.raw};
    \node[f] at (52,14.5) {kernel8-xv6\_wifi.img（新）};
    \node[f] at (52,11) {（早期）FS.IMG};
    \node[lab, anchor=east] at (96.5,5) {运行时读写};
    % readers
    \node[gbox, text width=36mm, minimum height=16mm] (gpu) at (20,61) {VideoCore GPU\\ROM → \code{bootcode.bin}\\→ \code{start.elf}\\（ARM 尚未运行）};
    \node[hbox, text width=36mm, minimum height=14mm] (drv) at (20,37) {xv6 内核驱动\\\code{fat32openroot()}\\\code{fat32readnext()}};
    \node[box, text width=36mm, minimum height=14mm] (usr) at (20,12) {用户程序经 \code{/boot}\\\code{ls cat mv tftpclient}\\\code{camshot}};
    \draw[rd] (51,61) -- node[lab, above]{读} (gpu.east);
    \draw[rd] (51,37) -- node[lab, above]{读} (drv.east);
    \draw[wr] (51,12) -- node[lab, above]{读写} (usr.east);
    \node[lab, align=left, anchor=west] at (102,61) {把 \code{kernel8-xv6\_wifi.img}\\装到物理 \code{0x80000}，\\然后放开 ARM};
    \node[lab, align=left, anchor=west] at (102,37) {\file{brcmfmac.c}、\file{mt7601u.c}\\每次上电下载给网卡};
    \node[lab, align=left, anchor=west] at (102,12) {VFS 挂载 \code{/boot}，\\按 \code{/etc/fstab} 读写};
  \end{tikzpicture}
  \caption{启动分区里的文件和它们的三类读者。同一个 FAT32 分区，先被 GPU 固件读（ARM 还没开始执行），再被 xv6 内核驱动读，最后作为 \code{/boot} 被用户程序读写}
  \label{fig:bootfs}
\end{figure}

\paragraph{第一组由 GPU 读，在 ARM 执行第一条指令之前.} 树莓派上电后先运行的是 VideoCore GPU：片内 ROM 读 \file{bootcode.bin}，它再装入 \file{start.elf} 和 \file{fixup.dat}。\file{start.elf} 解析 \file{config.txt}：

\begin{txtcode}
arm_64bit=1                   # 以 AArch64 启动 ARM 核
enable_uart=1                 # 打开 Mini UART，固定 core 时钟
uart_2ndstage=1               # 固件自己的启动日志也从串口输出
disable_commandline_tags=1    # 不往内存里写 ATAG/设备树
core_freq=250                 # VPU 时钟 250 MHz（Unicam 的下限，串口波特率也依赖它）
core_freq_min=250
kernel=kernel8-xv6_wifi.img   # 内核镜像文件名
camera_auto_detect=0          # 不让固件占用摄像头
\end{txtcode}

然后把 \file{kernel8-xv6\_wifi.img} 原样拷到物理地址 \code{0x80000}，才放开 ARM 核（第~\ref{chap:boot} 章）。所以内核镜像必须是裸二进制而不是 ELF，固件也从不理会 xv6 的根文件系统。

\paragraph{第二组由 xv6 内核驱动读.} 两块 Wi-Fi 网卡每次上电都要由主机下载固件：BCM43430/43455 需要 \file{.BIN}（R4 的执行代码）、\file{.TXT}（NVRAM）和 \file{.CLM}（法规数据库），MT7601U 需要 \file{MT7601U.BIN}。驱动在 \code{main()} 里初始化，这时还没有任何进程，更谈不上挂载点，所以它们直接调用 \code{fat32openroot()} 按 8.3 短名在根目录里找文件，再用 \code{fat32readnext()} 顺序读出。这也是挂起 bootfs 的 \code{rootdev\_init()} 必须排在 \code{brcmfmac\_driver\_init()} 之前的原因。

\paragraph{第三组由用户程序经 \code{/boot} 读写.} 进入用户态以后，\code{init} 按 \file{/etc/fstab} 把同一个分区挂到 \code{/boot}（\ref{sec:vfs} 节）。于是 \code{ls /boot}、\code{cat /boot/config.txt} 都走普通的 \code{open/read}；\code{tftpclient} 把新内核下载成 \file{/boot/kernel8-xv6\_wifi.img}，重启即生效；\code{camshot} 把超过 xv6 单文件上限的照片写成 \file{/boot/camera.bmp}，关机后插到电脑上就能看到。

\subsection{磁盘上的 FAT32}

FAT32 的结构可以用一句话概括：\emph{文件内容按簇存放，簇与簇之间的先后关系记在一张表（FAT）里，目录本身也是一个文件}。图~\ref{fig:fat-layout} 是它在 SD 卡上的位置和内部划分。

\begin{figure}[htbp]
  \centering
  \begin{tikzpicture}[x=1mm,y=1mm, font=\scriptsize]
    \tikzset{seg/.style={draw=structurecolor!70!black, minimum height=10mm, inner sep=1pt, align=center, anchor=west}}
    % whole card
    \node[lab, anchor=east] at (-1,30) {整张卡};
    \node[seg, fill=gray!20, minimum width=12mm] (mbr) at (0,30) {MBR\\扇区 0};
    \node[seg, fill=gray!5, minimum width=8mm] (gap) at (mbr.east) {空};
    \node[seg, fill=structurecolor!20, minimum width=40mm] (p1) at (gap.east) {p1 BOOTFS（FAT32）\\LBA 2048 起};
    \node[seg, fill=orange!15, minimum width=34mm] (p2) at (p1.east) {p2 xv6 原始分区\\LBA 208896 起};
    \node[seg, fill=gray!10, minimum width=26mm] (p3) at (p2.east) {p3 ext2 …};
    % zoom p1
    \node[lab, anchor=east] at (-1,8) {p1 内部};
    \node[seg, fill=structurecolor!30, minimum width=10mm] (bpb) at (0,8) {BPB\\扇区 0};
    \node[seg, fill=structurecolor!20, minimum width=10mm] (fsi) at (bpb.east) {FSInfo\\扇区 1};
    \node[seg, fill=structurecolor!10, minimum width=12mm] (rsv) at (fsi.east) {保留区\\其余扇区};
    \node[seg, fill=orange!25, minimum width=20mm] (fat1) at (rsv.east) {FAT \#1};
    \node[seg, fill=orange!25, minimum width=20mm] (fat2) at (fat1.east) {FAT \#2（副本）};
    \node[seg, fill=gray!8, minimum width=48mm] (data) at (fat2.east) {数据区：簇 2、3、4 …\\根目录也是一条簇链};
    \draw[gray!60, dashed] (p1.south west) -- (bpb.north west);
    \draw[gray!60, dashed] (p1.south east) -- (data.north east);
    % labels below
    \draw[decorate, decoration={brace, amplitude=3pt, mirror}] ([yshift=-1mm]bpb.south west) -- node[lab, below=3pt]{\code{reserved} 个扇区} ([yshift=-1mm]rsv.south east);
    \draw[decorate, decoration={brace, amplitude=3pt, mirror}] ([yshift=-1mm]fat1.south west) -- node[lab, below=3pt]{\code{fats} $\times$ \code{fat\_sectors}} ([yshift=-1mm]fat2.south east);
    \node[lab, anchor=north] at (fat1.south west |- 0,-4) {\code{fat\_lba}};
    \node[lab, anchor=north] at (data.south west |- 0,-4) {\code{data\_lba}};
    \draw[gray!60] (fat1.south west) -- ++(0,-6.5);
    \draw[gray!60] (data.south west) -- ++(0,-6.5);
  \end{tikzpicture}
  \caption{SD 卡与 FAT32 分区的磁盘布局。\code{fat32mount()} 从 BPB 读出各区的大小，算出 \code{fat\_lba} 和 \code{data\_lba}；簇号从 2 开始编号}
  \label{fig:fat-layout}
\end{figure}

\paragraph{MBR 和分区表.} 整张卡的第 0 扇区是 MBR：偏移 446 起是 4 个 16 字节的分区表项，最后两个字节是签名 \code{0x55 0xAA}。每个表项里第 4 字节是分区类型（\code{0x0B}/\code{0x0C} 是 FAT32），第 8、12 字节是起始 LBA 和扇区数。读分区表的是 \file{rootdev.c} 的 \code{rootdev\_init()}，而不是 FAT32 驱动。它先把四个表项全部打印出来，再分两轮挑选：第一轮找 block 1 带 xv6 superblock 魔数的分区做根文件系统（\code{setup\_raw\_root()}），第二轮把第一个 FAT32 分区的起点交给 \code{fat32mount()}。两轮互不依赖，所以根文件系统和启动分区可以各在各的分区里。FAT32 驱动的工作从 \code{fat32mount()} 拿到分区起点开始。

\paragraph{BPB：分区的第一个扇区.} FAT32 分区的第 0 扇区叫 BPB（BIOS Parameter Block），记录了整个分区的几何参数。\code{fat32mount()} 用到的字段如表~\ref{tab:bpb}。

\begin{table}[htbp]
  \centering
  \caption{\code{fat32mount()} 读取的 BPB 字段}
  \label{tab:bpb}
  \small
  \begin{tabularx}{\linewidth}{@{}llX@{}}
    \toprule
    偏移 & 字段 & 用途 \\
    \midrule
    11 & 每扇区字节数 & 必须是 512，否则拒绝 \\
    13 & 每簇扇区数 \code{spc} & 必须是 2 的幂；一簇通常 4\,KiB 或更大 \\
    14 & 保留扇区数 \code{reserved} & FAT 表从 \code{partition\_lba + reserved} 开始 \\
    16 & FAT 份数 \code{fats} & 通常为 2，第二份是备份，写入时要同步 \\
    32 & 总扇区数 & 与 FAT 大小一起算出总簇数 \\
    36 & 每份 FAT 的扇区数 & 数据区从 \code{fat\_lba + fats×fat\_sectors} 开始 \\
    44 & 根目录首簇 & 通常为 2；根目录也是一条簇链 \\
    48 & FSInfo 扇区号 & 内有“下一个空闲簇”的提示 \\
    510 & \code{0x55 0xAA} & 签名 \\
    \bottomrule
  \end{tabularx}
\end{table}

有了这些字段，任何簇号都能换算成 SD 卡上的绝对扇区号：

\begin{ccode}
fat_lba  = partition_lba + reserved;
data_lba = fat_lba + fats * fat_sectors;
static uint32 cluster_lba(uint32 cluster) {     // 簇号从 2 开始
  return diskmap.data_lba + (cluster - 2) * diskmap.sectors_per_cluster;
}
\end{ccode}

\code{fat32mount()} 对每个字段都做了合法性检查（2 的幂、不越界、不溢出）——SD 卡上的数据是外部输入，一个损坏的 BPB 不应该让内核算出一个指向别处的 LBA 然后去写它。

\subsection{FAT 表与簇链}

FAT 表是一个巨大的 32 位整数数组，下标就是簇号，第 $c$ 项的值是“簇 $c$ 之后的下一个簇”。只有低 28 位有效，高 4 位保留且写入时要保持原值；0 表示空闲，\code{0x0FFFFFF8} 及以上表示链尾（EOC）。一个 512 字节的扇区装 128 项，所以查第 $c$ 项只需读一个扇区：

\begin{ccode}
static uint32 fat_next(uint32 cluster) {
  uint32 sector = diskmap.fat_lba + cluster / 128;   // 128 = 512 / 4
  uint32 pos    = (cluster % 128) * 4;
  if(sdsector(sector, scratch, 0) < 0) panic("fat32: read FAT");
  return le32(scratch + pos) & 0x0fffffffU;
}
\end{ccode}

\code{fat\_set\_next()} 是对应的写操作，它要把同一个扇区在每一份 FAT 里都读改写一遍，保证两份副本一致。图~\ref{fig:fat-chain} 下半部分画的就是一条簇链：从目录项里拿到首簇，然后反复 \code{fat\_next()} 直到 EOC。

\subsection{目录项与长文件名}

\begin{figure}[htbp]
  \centering
  \begin{tikzpicture}[x=1mm,y=1mm, font=\scriptsize]
    \tikzset{cell/.style={draw=gray!60, minimum height=7mm, inner sep=1pt, align=center, anchor=west}}
    % directory entries
    \node[lab, anchor=west] at (0,66) {根目录簇中的连续 32 字节目录项};
    \node[cell, fill=structurecolor!10, minimum width=118mm] (l2) at (0,60) {LFN \#2（序号 \code{0x42}＝最后一片 | 第 2 片）：\code{"ifi.img"}，属性 \code{0x0F}，校验和 \code{c}};
    \node[cell, fill=structurecolor!10, minimum width=118mm] (l1) at (0,53) {LFN \#1（序号 \code{0x01}）：\code{"kernel8-xv6\_w"}，前 13 个 UTF-16 字符，属性 \code{0x0F}，校验和 \code{c}};
    \node[cell, fill=orange!15, minimum width=22mm] (sn) at (0,45) {\code{KERNEL\~{}1IMG}\\名字 [0:11)};
    \node[cell, fill=orange!8, minimum width=14mm] (at) at (sn.east) {属性\\{}[11]＝\code{0x20}};
    \node[cell, fill=gray!5, minimum width=30mm] (mi) at (at.east) {时间、日期等};
    \node[cell, fill=orange!25, minimum width=22mm] (hi) at (mi.east) {首簇高 16 位\\{}[20:22)};
    \node[cell, fill=orange!25, minimum width=16mm] (lo) at (hi.east) {首簇低 16 位\\{}[26:28)};
    \node[cell, fill=orange!8, minimum width=14mm] (sz) at (lo.east) {大小\\{}[28:32)};
    \node[lab, anchor=west] at (119,56.5) {LFN 片倒序存放};
    \node[lab, anchor=west] at (119,45) {短名目录项};
    % FAT table
    \node[lab, anchor=west, align=left] at (109,22) {FAT 表\\每项 4 字节，低 28 位有效};
    \foreach \i/\v/\c in {2/EOC/gray!8, 3/4/gray!8, 4/EOC/gray!8, 5/6/orange!25, 6/9/orange!25, 7/0/gray!3, 8/0/gray!3, 9/EOC/orange!25, 10/0/gray!3}{
      \node[cell, minimum width=11mm, fill=\c] (fat\i) at ({(\i-2)*12},22) {[\i]\enspace\code{\v}};
    }
    % data clusters
    
    \foreach \i/\c in {2/gray!8, 3/gray!8, 4/gray!8, 5/structurecolor!20, 6/structurecolor!20, 7/gray!3, 8/gray!3, 9/structurecolor!20, 10/gray!3}{
      \node[cell, minimum width=11mm, fill=\c] (dc\i) at ({(\i-2)*12},6) {簇 \i};
    }
    \node[lab, anchor=west, align=left] at (109,6) {数据区的簇\\簇 2 是根目录};
    % arrows
    \draw[-{Stealth[length=2mm]}, thick, orange!80!black, rounded corners] (lo.south) -- (lo.south |- 0,37.5) -- node[lab, below, pos=0.5]{首簇号（高 16 位拼低 16 位）＝ 5} (fat5.north |- 0,37.5) -- (fat5.north);
    \draw[-{Stealth[length=1.6mm]}, thick, orange!80!black] (fat5.south east) to[out=-30, in=-150] (fat6.south west);
    \draw[-{Stealth[length=1.6mm]}, thick, orange!80!black] (fat6.south east) to[out=-25, in=-155] (fat9.south west);
    \foreach \i in {5,6,9}{ \draw[gray!60, dashed] (fat\i.south) ++(0,-0.2) -- (dc\i.north); }
  \end{tikzpicture}
  \caption{长文件名、目录项与簇链。文件 \code{kernel8-xv6\_wifi.img} 占两片 LFN 加一个短名目录项；短名项里的首簇号是 5，沿 FAT 表 5→6→9 直到 EOC 就是它的全部数据簇}
  \label{fig:fat-chain}
\end{figure}

目录的内容是一串 32 字节的目录项（图~\ref{fig:fat-chain} 上半部分）。短名目录项里有 8.3 格式的名字（11 字节，大写、空格填充）、属性字节、首簇号（拆成高低两个 16 位字段，这是 FAT16 向 FAT32 扩展时留下的痕迹）和文件大小。扫描目录时，首字节有两个特殊值：\code{0x00} 表示后面再没有目录项，\code{0xE5} 表示这一项已被删除。属性字节的 \code{0x10} 表示子目录，\code{0x08} 表示卷标。

长文件名（LFN）是后来加的：属性等于 \code{0x0F} 的目录项不是文件，而是一片名字，每片装 13 个 UTF-16 字符，倒序排在它所属的短名目录项前面。第一片（物理上最靠前的）序号带 \code{0x40} 标志，表示“最后一片”；每片都存着短名的校验和，以便发现“LFN 和短名不配套”（例如老系统删了短名但留下了 LFN）。\code{fat\_find\_path()} 扫描目录时把遇到的 LFN 片按序号拼进一个缓冲区，遇到短名项时检查三件事——片序号连续到 1、校验和与短名一致、拼出的名字与要找的名字相同（不区分大小写）——满足才算长名匹配；否则再按短名比较。

创建文件时反过来：\code{fat32createfile()} 先用 \code{fat\_make\_short\_alias()} 造一个 \code{KERNEL\~{}1.IMG} 这样的短名（取前 6 个字母数字、加 \code{\~{}1}～\code{\~{}9} 避免重名），算出需要几片 LFN，用 \code{fat\_reserve\_root\_slots()} 在根目录里找到足够的连续空位（不够就给根目录分配新簇），按倒序写入 LFN 片和短名项，最后补一个全 0 的目录项作为结束标记。

本项目的 FAT32 驱动只处理\emph{根目录}：查找、创建、枚举、重命名都在根目录簇链里进行，不支持子目录，LFN 也只支持 ASCII 字符。对启动分区来说这已经足够——上面列出的所有文件都在根目录里。

\subsection{函数一览与调用关系}

\begin{figure}[htbp]
  \centering
  \begin{tikzpicture}[x=1mm,y=1mm, font=\scriptsize,
      fn/.style={box, font=\scriptsize, minimum height=6.5mm, inner sep=1.5pt},
      call/.style={-{Stealth[length=1.6mm]}, semithick, draw=structurecolor!70!black}]
    \foreach \y/\t in {64/{调用者}, 48/{\file{fat32.c} 接口}, 30/{目录与 FAT}, 14/{扇区}, 2/{设备}}{
      \fill[gray!5] (0,\y-6.5) rectangle (150,\y+6.5);
      \node[lab, align=center, text width=14mm] at (7,\y) {\t};
    }
    % callers
    \node[fn, text width=34mm] (c1) at (34,64) {Wi-Fi 驱动下载固件\\\file{brcmfmac.c}、\file{mt7601u.c}};
    \node[fn, text width=34mm] (c2) at (80,64) {VFS \code{/boot}：\code{fat\_vread}\\\code{fat\_vwrite}、\code{fat\_vreaddir}…};
    \node[fn, text width=30mm] (c3) at (128,64) {根块设备 \file{rootdev.c}\\（仅 \code{FS.IMG} 兼容模式）};
    % interface
    \node[fn, text width=30mm] (i1) at (30,48) {\code{fat32openroot}\\\code{fat32readerinit/readnext}};
    \node[fn, text width=40mm] (i2) at (80,48) {\code{fat32statpath} \code{fat32pread}\\\code{fat32createfile} \code{fat32writefile}\\\code{fat32truncatefile} \code{fat32renamefile}};
    \node[fn, text width=30mm] (i3) at (128,48) {\code{fat32mapfile}\\\code{fat32mapsector}};
    % helpers
    \node[fn, text width=37mm] (h1) at (36,30) {\code{fat\_find\_path}/\code{fat\_find\_root}\\\code{fat\_reserve\_root\_slots}\\\code{fat\_make\_lfn\_entry}};
    \node[fn, text width=40mm] (h2) at (86,30) {\code{fat\_next} \code{fat\_set\_next}\\\code{fat\_alloc\_cluster} \code{fat\_free\_chain}\\\code{fat\_mark\_fsinfo\_unknown}};
    % sector
    \node[fn, text width=70mm] (s1) at (64,14) {\code{cluster\_lba(c) = data\_lba + (c-2)×spc}\\全部经 \code{sdsector(lba, scratch, write)}，持 \code{fat32\_lock}};
    \node[fn, text width=70mm, fill=gray!10, draw=gray!60] (d1) at (64,2) {\file{sdhost.c}：CMD17/CMD24，PIO 搬运 512 字节};
    \draw[call] (c1) -- (i1); \draw[call] (c2) -- (i2); \draw[call] (c3) -- (i3);
    \draw[call] (i1) -- (h1); \draw[call] (i2) -- (h1); \draw[call] (i2) -- (h2); \draw[call] (i1.south east) -- (h2.north west);
    \draw[call] (h1) -- (s1); \draw[call] (h2) -- (s1);
    \draw[call] (i3.east) -- ++(3,0) |- (s1.east);
    \draw[call] (s1) -- (d1);
    \node[fn, text width=30mm, fill=gray!10, draw=gray!60] (boot) at (127,30) {启动时 \code{rootdev\_init}\\→ \code{fat32mount}};
  \end{tikzpicture}
  \caption{FAT32 驱动的调用关系。三类调用者共用同一组目录与 FAT 操作，最终都落到一个 512 字节的 \code{scratch} 扇区缓冲和 \code{sdsector()} 上}
  \label{fig:fat-calls}
\end{figure}

\file{fat32.c} 的函数可以按图~\ref{fig:fat-calls} 分成四层：

\begin{itemize}
  \item \textbf{挂载}：\code{fat32mount(lba, sectors)} 解析 BPB、读 FSInfo 里的空闲簇提示；\code{fat32\_is\_boot\_sector()} 判断一个扇区是不是 FAT32 引导扇区，\file{rootdev.c} 用它识别“没有分区表、整张卡就是一个 FAT32”的 superfloppy 格式；\code{fat32init()} 只初始化 \code{fat32\_lock}。分区表扫描和 xv6 根分区识别都不在这个文件里。
  \item \textbf{对外接口}：只读的 \code{fat32openroot()}、\code{fat32openpath()}、\code{fat32statpath()}、\code{fat32readdirroot()}、\code{fat32pread()}、\code{fat32readerinit()}/\code{fat32readnext()}；写入的 \code{fat32createfile()}、\code{fat32openwrite()}、\code{fat32writefile()}、\code{fat32truncatefile()}、\code{fat32renamefile()}。VFS 的 \code{fat32\_ops} 把它们包装成 \code{/boot} 上的 vnode 操作。
  \item \textbf{目录与 FAT}：上一小节的查找和 LFN 函数，以及 \code{fat\_next()}、\code{fat\_set\_next()}、\code{fat\_alloc\_cluster()}、\code{fat\_free\_chain()}、\code{fat\_mark\_fsinfo\_unknown()}。
  \item \textbf{文件扇区映射}：\code{fat32mapfile()} 记下根目录中某个文件的整条簇链，\code{fat32mapsector()} 把“文件内第 $s$ 个扇区”换成卡上的扇区号。它们只为兼容早期把根镜像存成 \code{FS.IMG} 文件的旧卡而保留，调用者是 \file{rootdev.c}。FAT32 只回答“这个扇区在哪里”，从不解释扇区里的 xv6 数据。
\end{itemize}

两个读接口的区别值得一提。\code{fat32pread(file, offset, …)} 是无状态的：每次调用都从首簇沿链走 $\lfloor\text{offset}/\text{簇大小}\rfloor$ 步，适合 VFS 偶尔的随机读。如果用它按 512 字节分块读一个 600\,KiB 的 Wi-Fi 固件，总共要走 $O(n^2)$ 次 \code{fat\_next()}。\code{fat32reader} 则记住当前簇和簇内偏移，顺序读时每跨一个簇只查一次 FAT——固件下载就用它。

\subsection{写入：分配、追加与截断}

\code{fat32writefile()} 只支持追加：要求写入位置正好等于文件当前大小。写入过程是：

\begin{enumerate}
  \item 重新按路径查找目录项，核对首簇和大小与调用者手里的 \code{fat32\_file} 一致，防止文件已被别人改过；
  \item 写入跨进新簇时调用 \code{fat\_alloc\_cluster()}：从 FSInfo 提示的位置开始找 FAT 值为 0 的项，把它标成 EOC（每份 FAT 都写），把这一簇清零；再用 \code{fat\_set\_next()} 把它挂到链尾；
  \item 按扇区写数据：整扇区直接写，不满一个扇区的先读出来再改再写回；
  \item 最后一次性更新目录项里的首簇和大小。
\end{enumerate}

\code{fat32truncatefile()}（\code{O\_TRUNC}）的顺序恰好相反：先把目录项的首簇和大小清零并写回，再释放旧簇链。这个顺序的意义在于断电时的后果：先改目录项，最坏情况是旧簇成了没人引用的“丢失簇”，浪费一点空间；如果先释放簇链，断电后目录项会指向已被标为空闲、随时可能分给别人的簇，那才是真正的损坏。分配和释放之后，\code{fat\_mark\_fsinfo\_unknown()} 把 FSInfo 里的空闲簇计数标成“未知”（\code{0xFFFFFFFF}），只更新“下一个空闲簇”提示，避免留下一个过时的计数让别的系统误信。

两处性能问题都是真机上暴露出来的，修法相同——缓存当前 FAT 扇区：
\begin{itemize}
  \item \code{fat\_alloc\_cluster()} 最初对每个候选簇调用一次 \code{fat\_next()}，同一个扇区要重复读 128 次，根目录一长就像卡死了一样；现在扫描时只在跨扇区时才读；
  \item \code{fat\_free\_chain()} 最初对每个簇做一次读和两份 FAT 的读改写，截断一张 320$\times$240 的 BMP 要五百多次 SD 事务（第~\ref{chap:drivers} 章 \ref{sec:cam-bugs} 小节）；现在在一个扇区内批量清除链项，换扇区时才把整扇区写到每份 FAT。
\end{itemize}

\paragraph{并发与局限.} 所有接口都持同一把自旋锁 \code{fat32\_lock}，共用一个 512 字节的 \code{scratch} 扇区缓冲，没有块缓存，也没有日志。路径解析、LFN 字符、目录槽这些较大的临时数组也放在一个静态工作区 \code{fat\_work} 里，而不是放在栈上——每个进程的内核栈只有 4\,KiB，它们曾经把创建文件的调用链撑爆（第~\ref{chap:drivers} 章 \ref{sec:cam-bugs} 小节）；既然有 \code{fat32\_lock} 串行化，一个工作区就够用。这让实现很简单，代价也很清楚：持锁期间 \code{sdsector()} 以 PIO 方式等待 SD 卡，一次大文件写会长时间占着这把锁；同一时间只能有一个写者；写入途中断电可能损坏分区。启动分区平时几乎只读，偶尔更新内核或写一张照片，这些代价是可以接受的，但修改它之前最好备份。

\section{xv6 根文件系统：superblock、inode 与目录}
\label{sec:xv6fs}

前一节分析了启动分区上的 FAT32，这一节用同样的方法看 p2 上的 xv6 根文件系统：它怎么被认出来，磁盘上每一块放的是什么，文件和目录怎样表示，一个路径怎样一步步变成 SD 卡上的扇区。\ref{sec:rootfs} 节讲过根文件系统可以放在哪里；这里只讨论现在的“原始分区”方式。

\subsection{根文件系统怎样被认出来}

\code{rootdev\_init()} 扫描 MBR 时，对每个既不是 FAT32 也不是扩展分区的主分区调用 \code{setup\_raw\_root()}（两者都在 \file{kernel/rootdev.c}）：

\begin{ccode}
static int setup_raw_root(uint32 start, uint32 sectors) {
  // xv6 的第 1 块是 superblock；块大小 1024，所以它在分区起点之后第 2 个扇区
  if(start == 0 || sectors < FSSIZE * SECTORS_PER_BLOCK ||
     sdsector(start + SECTORS_PER_BLOCK, sector, 0) < 0 ||
     le32(sector) != FSMAGIC || le32(sector + 4) != FSSIZE)
    return -1;
  root.mode = ROOT_RAW;
  root.lba = start;
  root.sectors = sectors;
  return 0;
}
\end{ccode}

判断只看两样东西：superblock 开头的魔数 \code{0x10203040}，以及紧跟着的 \code{size} 是否等于内核编译时的 \code{FSSIZE}。MBR 里的分区类型不作数——macOS 的 \code{diskutil} 只能把这个分区建成 Linux 的 \code{0x83} 类型，和真正的 ext2 分区一样，靠魔数才能分开。\code{size} 的检查同样重要：如果有人改了 \code{FSSIZE} 重新编译内核却没重新烧 \code{fs.img}，这里会拒绝，而不是让新内核去解释旧格式的磁盘。识别成功后启动日志打印：

\begin{txtcode}
rootdev: raw xv6 root p2 lba=208896 sectors=819200 size=400 MiB
\end{txtcode}

此后块缓存的每一次读写都经过 \code{rootdev\_rw()}：

\begin{ccode}
void rootdev_rw(struct buf *b, int write) {         // kernel/rootdev.c
  uint32 first = b->blockno * SECTORS_PER_BLOCK;    // 1 块 = 2 扇区
  gate_enter();                                     // /dev/sdroot 更新时在此挡住
  for(int i = 0; i < SECTORS_PER_BLOCK; i++)
    if(sdsector(root_sector_lba(first + i),
                b->data + i * SECTOR_SIZE, write) < 0)
      panic(write ? "rootdev: write block" : "rootdev: read block");
  gate_exit();
}
\end{ccode}

原始分区模式下 \code{root\_sector\_lba()} 只做一次加法：扇区号 $=$ \code{root.lba} $+$ 块号 $\times 2 + i$。整条路径不经过 \file{fat32.c} 的任何函数。分区有 400\,MiB，文件系统只用前 32\,MiB；可分配多少块由 superblock 决定，与分区大小无关。

\subsection{superblock 与磁盘布局}

\begin{figure}[htbp]
  \centering
  \begin{tikzpicture}[x=1mm,y=1mm, font=\scriptsize]
    \tikzset{seg/.style={draw=structurecolor!70!black, minimum height=11mm, inner sep=1pt, align=center, anchor=west}}
    \node[seg, fill=gray!15, minimum width=10mm] (b0) at (0,10) {启动块\\0};
    \node[seg, fill=orange!30, minimum width=13mm] (sb) at (b0.east) {superblock\\1};
    \node[seg, fill=structurecolor!25, minimum width=24mm] (lg) at (sb.east) {日志\\2～31（30 块）};
    \node[seg, fill=structurecolor!15, minimum width=26mm] (in) at (lg.east) {inode 区\\32～44（13 块）};
    \node[seg, fill=orange!15, minimum width=20mm] (bm) at (in.east) {位图\\45～49（5 块）};
    \node[seg, fill=gray!5, minimum width=45mm] (dt) at (bm.east) {数据区\\50～32767（32718 块）};
    % zooms
    \node[lab, anchor=north, align=center] at (lg.south) {第一块是日志头：\\\code{n} + 30 个块号};
    \node[lab, anchor=north, align=center] at (in.south) {每块 16 个 \code{dinode}\\共 200 个 inode};
    \node[lab, anchor=north, align=center] at (bm.south) {每块 8192 位\\1 位管 1 块};
    \node[lab, anchor=north, align=center] at (dt.south) {文件内容、目录内容、\\间接块};
    \draw[-{Stealth[length=1.6mm]}, orange!70!black] (sb.north) to[out=70, in=110] node[lab, above]{\code{logstart}} (lg.north);
    \draw[-{Stealth[length=1.6mm]}, orange!70!black] (sb.north) to[out=60, in=115] node[lab, above, pos=0.6]{\code{inodestart}} (in.north);
    \draw[-{Stealth[length=1.6mm]}, orange!70!black] (sb.north) to[out=55, in=120] node[lab, above, pos=0.7]{\code{bmapstart}} (bm.north);
    \node[lab, anchor=south west] at (0,36) {块号（1\,KiB 一块）；分区内的扇区号 $=$ 块号 $\times 2$};
  \end{tikzpicture}
  \caption{\code{mkfs} 生成的 32\,MiB 根文件系统布局。所有数字都由 \code{FSSIZE}、\code{NINODES}、\code{LOGSIZE} 算出，并写进第 1 块的 superblock}
  \label{fig:xv6-layout}
\end{figure}

superblock 是第 1 块（第 0 块是从 Unix 继承下来的“启动块”，xv6 不用），内容就是 8 个 32 位整数：

\begin{table}[htbp]
  \centering
  \caption{\code{fs.img} 的 superblock（\code{struct superblock}，\file{kernel/fs.h}）}
  \label{tab:xv6-sb}
  \small
  \begin{tabularx}{\linewidth}{@{}llX@{}}
    \toprule
    字段 & 值 & 含义与来源 \\
    \midrule
    \code{magic} & \code{0x10203040} & \code{FSMAGIC}；\code{fsinit()} 和 \code{setup\_raw\_root()} 都检查它 \\
    \code{size} & 32768 & 总块数，\code{FSSIZE}；32\,MiB \\
    \code{nblocks} & 32718 & 数据块数 $=$ \code{size} $-$ 元数据块数 50 \\
    \code{ninodes} & 200 & \code{mkfs} 里的 \code{NINODES} \\
    \code{nlog} & 30 & \code{LOGSIZE} $=$ \code{MAXOPBLOCKS}$\times 3$ \\
    \code{logstart} & 2 & 紧跟 superblock \\
    \code{inodestart} & 32 & $2 + 30$；inode 区 $200/16+1=13$ 块 \\
    \code{bmapstart} & 45 & $32 + 13$；位图 $32768/8192+1=5$ 块 \\
    \bottomrule
  \end{tabularx}
\end{table}

这些数都是 \code{mkfs} 在宿主机上算出来写进去的（运行 \code{make fs.img} 时会打印 \code{nmeta 50 (boot, super, log blocks 30 inode blocks 13, bitmap blocks 5) blocks 32718 total 32768}）。内核启动时 \code{fsinit()} 读第 1 块、核对魔数、交给日志层，之后所有“某个东西在哪一块”的问题都由 superblock 加两个宏回答：

\begin{ccode}
#define IPB           (BSIZE / sizeof(struct dinode))   // 每块 16 个 inode
#define IBLOCK(i, sb) ((i) / IPB + sb.inodestart)       // inode i 所在的块
#define BPB           (BSIZE * 8)                       // 每个位图块管 8192 块
#define BBLOCK(b, sb) ((b) / BPB + sb.bmapstart)        // 块 b 的位所在的位图块
\end{ccode}

\paragraph{日志区.} 30 块中第一块是日志头 \code{struct logheader}：一个计数 \code{n} 加上 \code{n} 个“这些块属于哪个目标块号”。其余是被修改的块的副本。提交时先写副本、再写日志头（这一写就是提交点）、再把副本拷回原位、最后把 \code{n} 清零；\code{fsinit()} → \code{initlog()} 在启动时检查日志头，\code{n} 不为 0 就重做一遍。

\paragraph{位图.} 每一位对应一块，1 表示已用。\code{mkfs} 把前 \code{freeblock} 块（元数据加上镜像里已经写入的全部文件内容）一次性标成已用；运行时 \code{balloc()} 从头扫描位图找第一个 0 位，置 1、\code{log\_write}、再把新块清零。早期 \code{FSSIZE} 只有 1000 块时，光是打包进镜像的程序就用掉 986 块，首次启动创建 \code{/etc} 配置时 \code{balloc} 找不到 0 位，panic “out of blocks”——这就是 \code{FSSIZE} 改成 32768 的原因。

\subsection{inode：文件的元数据与块映射}

\begin{figure}[htbp]
  \centering
  \begin{tikzpicture}[x=1mm,y=1mm, font=\scriptsize]
    \tikzset{f/.style={draw=gray!60, minimum height=6mm, inner sep=1pt, align=center, anchor=west}}
    \node[lab, anchor=west] at (0,47) {\code{struct dinode}（64 字节，inode 区每块 16 个）};
    \node[f, fill=orange!20, minimum width=12mm] (t) at (0,40) {\code{type}\\2 B};
    \node[f, fill=gray!8, minimum width=12mm] (ma) at (t.east) {\code{major}\\2 B};
    \node[f, fill=gray!8, minimum width=12mm] (mi) at (ma.east) {\code{minor}\\2 B};
    \node[f, fill=orange!10, minimum width=12mm] (nl) at (mi.east) {\code{nlink}\\2 B};
    \node[f, fill=orange!10, minimum width=12mm] (sz) at (nl.east) {\code{size}\\4 B};
    \node[f, fill=structurecolor!15, minimum width=60mm] (ad) at (sz.east) {\code{addrs[0..11]}：12 个直接块号\\48 B};
    \node[f, fill=structurecolor!30, minimum width=16mm] (ind) at (ad.east) {\code{addrs[12]}\\间接块号};
    % data blocks
    \foreach \i in {0,...,4}{
      \draw[fill=gray!5, draw=gray!60] (62+\i*7,22) rectangle ++(6,6);
    }
    \node[lab] at (98,25) {…};
    \node[lab, anchor=west] at (62,18) {直接块：文件第 0～11 KiB};
    \foreach \i in {0,...,3}{ \draw[-{Stealth[length=1.2mm]}, gray!60] (64+\i*10,37) -- (65+\i*7,28); }
    \draw[fill=structurecolor!20, draw=structurecolor!70!black] (122,18) rectangle ++(18,10) node[midway, align=center]{间接块\\256 个块号};
    \draw[-{Stealth[length=1.6mm]}, thick, structurecolor!70!black] (ind.south) -- (131,28);
    \foreach \i in {0,...,4}{
      \draw[fill=gray!5, draw=gray!60] (110+\i*7,2) rectangle ++(6,6);
    }
    \node[lab] at (147,5) {…};
    \draw[-{Stealth[length=1.6mm]}, gray!60] (131,18) -- (124,8);
    \draw[-{Stealth[length=1.6mm]}, gray!60] (131,18) -- (138,8);
    \node[lab, anchor=east] at (108,5) {文件第 12～267 KiB};
    \node[lab, anchor=west, align=left] at (0,24) {\code{type}：0 空闲，1 目录，\\2 普通文件，3 设备\\[2pt]\code{nlink}：有几个目录项指向它\\[2pt]最大文件 $12+256=268$ 块};
  \end{tikzpicture}
  \caption{磁盘上的 inode 和它的块映射。\code{bmap(ip, bn)} 把文件内第 \code{bn} 块翻译成磁盘块号：小于 12 直接查 \code{addrs[bn]}，否则读间接块再查第 \code{bn-12} 项；遇到 0 就当场 \code{balloc()} 一块}
  \label{fig:xv6-inode}
\end{figure}

每个文件、目录和设备节点都由一个 64 字节的磁盘 inode 描述（图~\ref{fig:xv6-inode}）。它不包含文件名——名字在目录里——只有类型、链接数、大小和数据块的位置。内存里的 \code{struct inode} 在此之外再加上 \code{dev}、\code{inum}、引用计数 \code{ref}、睡眠锁和 \code{valid} 标志（是否已从磁盘读入）。\code{iget()} 只在 inode 缓存里占一个槽、不读盘；\code{ilock()} 第一次加锁时才 \code{bread(IBLOCK(...))} 把磁盘 inode 拷进来；修改后 \code{iupdate()} 写回。

块映射由 \code{bmap()} 完成：

\begin{ccode}
static uint bmap(struct inode *ip, uint bn) {
  uint addr, *a;
  struct buf *bp;
  if(bn < NDIRECT){                                   // 前 12 块：直接块
    if((addr = ip->addrs[bn]) == 0)
      ip->addrs[bn] = addr = balloc(ip->dev);
    return addr;
  }
  bn -= NDIRECT;
  if(bn < NINDIRECT){                                 // 之后 256 块：经间接块
    if((addr = ip->addrs[NDIRECT]) == 0)
      ip->addrs[NDIRECT] = addr = balloc(ip->dev);
    bp = bread(ip->dev, addr);
    a = (uint*)bp->data;
    if((addr = a[bn]) == 0){
      a[bn] = addr = balloc(ip->dev);
      log_write(bp);
    }
    brelse(bp);
    return addr;
  }
  panic("bmap: out of range");
}
\end{ccode}

所以单个文件最大 $12+256=268$ 块，即 268\,KiB。这是 xv6 根文件系统最明显的限制：Wi-Fi 固件、640$\times$480 的 BMP（900\,KB）、原始帧（375\,KB）都放不下，只能放在 FAT32 的 \code{/boot}（\ref{sec:fat32} 节）。

根文件系统里几个特殊的 inode 号：1 号永远是根目录（\code{ROOTINO}），\code{namex()} 遇到以 \code{/} 开头的路径就从它开始；2 号是 \code{mkfs} 建的 \code{/bin}；之后是按 \code{Makefile} 顺序装入的每个用户程序。\code{init} 是唯一一个在根目录下也有名字的程序（启动代码执行的是 \code{exec("init")}），所以 \code{/init} 和 \code{/bin/init} 是同一个 inode 的两个目录项，\code{mkfs} 把它的 \code{nlink} 设为 2。

\subsection{目录与路径查找}

\begin{figure}[htbp]
  \centering
  \begin{tikzpicture}[x=1mm,y=1mm, font=\scriptsize,
      ino/.style={draw=orange!70!black, fill=orange!12, rounded corners=1pt, align=left, inner sep=2pt, text width=34mm, anchor=north west},
      de/.style={draw=gray!60, fill=gray!4, minimum height=4.6mm, minimum width=34mm, inner sep=1pt, anchor=north west, font=\scriptsize\ttfamily},
      hit/.style={de, fill=structurecolor!20, draw=structurecolor!70!black},
      go/.style={-{Stealth[length=1.8mm]}, thick, draw=structurecolor!80!black},
      lk/.style={-{Stealth[length=1.6mm]}, semithick, draw=orange!70!black}]
    % column 1: root
    \node[lab, anchor=west] at (0,62) {\ding{172} \code{iget(ROOTDEV, 1)}};
    \node[ino] (i1) at (0,58) {inode 1：\code{type=1} 目录\\\code{size=1024}，\code{addrs[0]=50}};
    \node[lab, anchor=west] at (0,44) {块 50（根目录内容）};
    \node[de] (r0) at (0,41) {\ \ 1  .};
    \node[de] (r1) at (r0.south west) {\ \ 1  ..};
    \node[hit] (r2) at (r1.south west) {\ \ 2  bin};
    \node[de] (r3) at (r2.south west) {\ \ k  init};
    \node[de] (r4) at (r3.south west) {\ \ …  dev etc boot…（init 建）};
    \node[de] (r5) at (r4.south west) {\ \ 0  （空槽）};
    \draw[lk] (i1.south) ++(-6,0) -- ++(0,-3.5);
    % column 2: /bin
    \node[lab, anchor=west] at (52,62) {\ding{173} \code{dirlookup(root,"bin")}};
    \node[ino] (i2) at (52,58) {inode 2：\code{type=1} 目录\\\code{size=}若干 KiB，\code{addrs[0]=51}…};
    \node[lab, anchor=west] at (52,44) {块 51…（\code{/bin} 的内容）};
    \node[de] (b0) at (52,41) {\ \ 2  .};
    \node[de] (b1) at (b0.south west) {\ \ 1  ..};
    \node[de] (b2) at (b1.south west) {\ \ 3  cat};
    \node[de] (b3) at (b2.south west) {\ \ …};
    \node[hit] (b4) at (b3.south west) {\ \ m  ls};
    \node[de] (b5) at (b4.south west) {\ \ …};
    \draw[lk] (i2.south) ++(-6,0) -- ++(0,-3.5);
    % column 3: ls
    \node[lab, anchor=west] at (104,62) {\ding{174} \code{dirlookup(bin,"ls")}};
    \node[ino] (i3) at (104,58) {inode m：\code{type=2} 文件\\\code{size}＝ELF 大小\\\code{addrs[0..11]}、\code{addrs[12]}};
    \node[draw=gray!60, fill=gray!5, align=center, text width=34mm, anchor=north west, inner sep=2pt] (dl) at (104,41) {数据块：ELF 头、\\程序段……\\（\code{exec} 用 \code{readi} 读取）};
    \draw[lk] (i3.south) ++(-6,0) -- ++(0,-6.5);
    % lookups
    \draw[go] (r2.east) -- ++(4,0) |- (i2.west);
    \draw[go] (b4.east) -- ++(4,0) |- (i3.west);
    \node[lab, anchor=west, align=left] at (0,6) {每个目录项 32 字节：\code{ushort inum} + \code{char name[30]}，每块 32 项；\code{inum=0} 表示空槽。\\\code{/init} 与 \code{/bin/init} 是同一个 inode k 的两个名字，所以它的 \code{nlink=2}。};
  \end{tikzpicture}
  \caption{\code{namei("/bin/ls")} 的查找过程。目录本身也是文件，内容是一串目录项；\code{namex()} 每次用 \code{skipelem()} 取出一个路径分量，在当前目录的内容里 \code{dirlookup()}，得到下一个 inode。块号 50、51 是 \code{mkfs} 的分配顺序决定的，k、m 取决于 \code{Makefile} 里程序的排列顺序}
  \label{fig:xv6-namei}
\end{figure}

目录就是 \code{type=1} 的文件，内容是一串 32 字节的 \code{struct dirent}：2 字节 inode 号加 30 字节文件名（\code{DIRSIZ} 从原版的 14 扩到 30，见 \ref{sec:layout} 节），一块恰好 32 项，\code{inum=0} 表示空槽。\code{mkfs} 把目录大小向上取整到整块。目录操作只有两个基本函数：\code{dirlookup(dp, name)} 顺序读目录内容、比较名字，命中就 \code{iget()} 对应的 inode；\code{dirlink(dp, name, inum)} 先确认不重名，再找第一个空槽（没有就追加在末尾）写入新目录项。

路径查找是 \code{namex()}，图~\ref{fig:xv6-namei} 画的是 \code{namei("/bin/ls")}：

\begin{ccode}
if(*path == '/') ip = iget(ROOTDEV, ROOTINO);        // 绝对路径从 inode 1 开始
else             ip = idup(myproc()->cwd);            // 相对路径从当前目录开始
while((path = skipelem(path, name)) != 0){           // 依次取出 "bin"、"ls"
  ilock(ip);
  if(ip->type != T_DIR){ iunlockput(ip); return 0; }
  if(nameiparent && *path == '\0'){ iunlock(ip); return ip; }   // 停在父目录
  if((next = dirlookup(ip, name, 0)) == 0){ iunlockput(ip); return 0; }
  iunlockput(ip);
  ip = next;
}
\end{ccode}

注意加锁的顺序：每一步只锁当前目录，查到下一级以后先 \code{iunlockput()} 当前目录再进入下一级，任何时刻最多持有一把 inode 锁，所以不会因为两个进程以相反顺序查找路径而死锁。\code{nameiparent()} 在最后一个分量之前停下，返回父目录并把最后一个名字留给调用者——\code{create}、\code{unlink}、\code{link} 都需要它。

根目录的内容分两部分：\code{mkfs} 写入的 \code{.}、\code{..}、\code{bin}、\code{init}；以及 \code{init} 首次启动时创建的 \code{dev}（含 \code{console}、\code{ttyS0}、\code{input/event0}、\code{video0}、\code{sdroot} 等设备节点）、\code{etc}（\code{fstab} 等配置）、\code{proc}、\code{boot}、\code{mnt/ext2}、\code{root}、\code{usr/bin}。其中 \code{proc}、\code{boot}、\code{mnt/ext2} 只是空目录，挂载以后会被 VFS 的挂载表“盖住”（\ref{sec:vfs} 节）。设备节点是 \code{type=3} 的 inode，不占数据块，只在 \code{major}/\code{minor} 里记下交给哪个驱动。

\subsection{从路径到 SD 扇区}

把这一节和前面几节串起来，读 \code{/bin/ls} 的第 0 字节要经过这些层：

\begin{txtcode}
namei("/bin/ls")             -> inode m（经 inode 1、inode 2 的目录内容）
readi(ip, ..., off=0, n)     -> bn = off / 1024 = 0
bmap(ip, 0)                  -> 块号 B = ip->addrs[0]
bread(ROOTDEV, B)            -> 块缓存未命中，调用 rootdev_rw(b, 0)
root_sector_lba(2B), (2B+1)  -> 扇区 208896 + 2B 和 208896 + 2B + 1
sdsector(lba, buf, 0)        -> SDHOST CMD17，PIO 读 512 字节
\end{txtcode}

一个 1\,KiB 的 xv6 块对应两个 512 字节扇区。写入的路径与此相同，只是中间多了日志：\code{writei()} 修改的块经 \code{log\_write()} 记入日志，\code{end\_op()} 提交时才真正经 \code{rootdev\_rw()} 写到 SD 卡。

\subsection{\code{mkfs} 怎样生成 \code{fs.img}}

\code{mkfs} 是一个在宿主机上运行的普通 C 程序，它和内核共用 \file{kernel/fs.h}，因此磁盘格式只有一份定义：

\begin{enumerate}
  \item 按 \code{FSSIZE}、\code{NINODES}、\code{LOGSIZE} 算出各区大小，把 32768 块全部写 0，再在第 1 块写入 superblock；
  \item \code{ialloc(T\_DIR)} 得到 inode 1，写入 \code{.} 和 \code{..}（都指向自己）；再建 inode 2 作为 \code{/bin}，在根目录里加 \code{bin} 项；
  \item 对命令行上的每个程序（\code{\_ls}、\code{\_cat}…去掉前缀下划线）分配一个 inode、在 \code{/bin} 里加目录项、用 \code{iappend()} 把文件内容逐块追加——它和内核的 \code{bmap()} 一样，前 12 块直接放、之后经间接块；\code{init} 额外在根目录里加一项；
  \item 把两个目录的大小取整到整块；最后 \code{balloc(freeblock)} 把已经用掉的前 \code{freeblock} 块在位图里标成 1。
\end{enumerate}

\code{mkfs} 只按顺序分配块和 inode，不需要日志，也不处理删除；所有整数都用 \code{xint()}/\code{xshort()} 写成小端序，与 AArch64 内核一致。

\subsection{在系统里更新根文件系统}

根文件系统的格式一旦改动（例如 \code{DIRSIZ} 从 14 改成 30），就必须整体重写 p2。最初只能把 SD 卡拔到电脑上 \code{dd}；后来内核增加了一个受保护的字符设备 \code{/dev/sdroot}，它只允许写启动时 \code{setup\_raw\_root()} 选中的那个分区，绝不暴露整张卡。更新流程是：

\begin{shcode}
/bin/tftpclient 192.168.0.195 fs.img          # 下载到 /boot（FAT32）
/bin/dd if=/boot/fs.img of=/dev/sdroot bs=32768
\end{shcode}

\code{dd} 在写之前先检查镜像大小恰好是 \code{FSSIZE}$\times$\code{BSIZE}、把整个镜像读进内存并校验 superblock 的魔数、大小和各区起点——因为一旦开始覆盖 p2，就再也不能从旧的根文件系统读任何东西，连 \code{dd} 自己的代码也不行（用户程序是静态链接的，已全部在内存里）。打开 \code{/dev/sdroot} 时内核冻结根文件系统：不再接受新的事务，等已有事务提交、读者退出；然后每写一个扇区就读回比较一次。结束时打印：

\begin{txtcode}
sdroot: verified 32/32 MiB
sdroot: update complete and verified; reboot now
dd: 33554432 bytes installed and read-back verified
dd: root filesystem is frozen; power-cycle Raspberry Pi now
\end{txtcode}

冻结是单向的，无论成败都必须断电重启——内存里的块缓存和 inode 缓存描述的是旧磁盘，不能再拿来解释新磁盘。

\section{烧录 SD 卡}

\subsection{分区布局}

一张为 xv6 准备的 SD 卡分成四个区（表~\ref{tab:sdparts}）。

\begin{table}[htbp]
  \centering
  \caption{SD 卡分区（63.8\,GB 卡的实例）}
  \label{tab:sdparts}
  \small
  \begin{tabular}{@{}lllrl@{}}
    \toprule
    分区 & 名称 & 格式 & 起始 LBA & 用途 \\
    \midrule
    p1 & BOOTFS & FAT32 & 2048 & 固件、\code{config.txt}、内核、Wi-Fi 固件 \\
    p2 & XV6FS & 原始 & 208896 & xv6 根文件系统（400\,MiB，用到 32\,MiB） \\
    p3 & EXT2 & ext2 & 1030144 & Linux 文件交换 \\
    p5 & SPFS & 原始 & 123654144 & 预留给其他教学文件系统 \\
    \bottomrule
  \end{tabular}
\end{table}

\subsection{步骤}

以 macOS 为例（\code{disk4} 必须先用 \code{diskutil list} 确认，写错磁盘会毁掉数据）：

\begin{shcode}
# 1. 卸载整张卡并重新分区
diskutil unmountDisk force /dev/disk4
diskutil partitionDisk /dev/disk4 4 MBR \
  FAT32 BOOTFS 100MiB  UFSD_EXTFS XV6FS 400MiB \
  UFSD_EXTFS EXT2 R    UFSD_EXTFS SPFS 500MiB

# 2. 写入 xv6 根文件系统（只能写分区 2，绝不能写整卡 /dev/rdisk4）
make fs.img
make install-rpi3-rawfs RPI3_XV6_DEV=/dev/rdisk4s2

# 3. 把分区 3 格式化成不带 ext4 特性的 ext2
sudo mke2fs -F -t ext2 -L EXT2 /dev/rdisk4s3

# 4. 在 BOOTFS 中放入树莓派官方固件（bootcode.bin、start.elf、fixup.dat 等），
#    然后安装内核、config.txt 和 Wi-Fi 固件
make install-rpi3 RPI3_XV6_DEV=/dev/rdisk4s2
\end{shcode}

\code{Makefile} 对 \code{RPI3\_XV6\_DEV} 做了检查：变量为空时只更新内核、配置和固件，并给出警告；不是 \code{/dev/diskNsM} 形式的分区设备则拒绝执行。这些防护是吃过亏之后加上的——\code{dd} 写错一个字母就会覆盖整张卡的分区表。

\subsection{日常更新}

改了内核之后不必每次拔卡：运行 \code{make}，内核会被拷到宿主机的 TFTP 目录；在树莓派上执行

\begin{shcode}
tftp 192.168.0.195 kernel8-xv6_wifi.img
\end{shcode}

文件就会下载到 \code{/boot}（FAT32 启动分区），重启后生效。只有用户程序或文件系统格式变化时，才需要重新烧录分区 2。
